\documentclass[11pt]{article}
\usepackage[margin=1in]{geometry}
\usepackage{amsmath,amssymb,amsthm,hyperref}
\newtheorem{proposition}{Proposition}
\newcommand{\R}{\mathbb R}
\newcommand{\Al}{A_{\bullet}}
\newcommand{\Au}{A^{\bullet}}
\title{Coinciding transition attractors need not attract at the threshold}
\author{Principia Math in Progress}
\date{12 September 2026}
\begin{document}
\maketitle
\paragraph{Task 1. Resolve MathDB 354469.}
Conjecture 6.1 of Le\'sniak, Snigireva, Strobin and Vince
\cite{source}, p.~5425, asks whether, under (H1)--(H3), equality
of a lower transition attractor with one upper transition attractor
forces uniqueness of the upper transition attractor and attraction by
the terminal IFS. The terminal-attraction assertion is false.

Here (H1) requires continuous dependence of each map on the parameter
at every point; (H2) requires a Lipschitz constant below one for each
parameter $t<1$; and (H3) requires the individual fixed points to have
limits as $t\uparrow1$. The source defines an attractor using
Hausdorff convergence from every nonempty compact subset of some open
neighborhood (p.~5398). Mere invariance is insufficient.

\begin{proposition}
On the complete metric space $\R$ with its usual distance, put
\[
 r_t=\frac{1+t}{2},\qquad
 F_t=\{g_t,f_t\},\qquad
 g_t(x)=r_tx,\quad f_t(x)=\frac{r_t}{2}x
 \quad(0\leq t\leq1).
\]
This family satisfies (H1)--(H3), has $\Al=\Au=\{0\}$ and a unique
upper transition attractor, but $\{0\}$ is not an attractor of $F_1$.
\end{proposition}

\begin{proof}
The two maps are distinct, invertible and continuously parameterized,
including at $t=0$. For $0\leq t<1$, their maximal Lipschitz constant
is $r_t<1$, and both have the unique fixed point zero. Thus (H1)--(H3)
hold and the set $Q$ of limiting fixed points is $\{0\}$.

For every nonempty compact $K\subset\R$ and integer $n\geq0$,
\begin{equation}\label{iterate}
 F_t^n(K)=\bigcup_{k=0}^{n} r_t^n2^{-k}K.
\end{equation}
Indeed, a word of length $n$ containing $k$ copies of $f_t$ multiplies
every point by $r_t^n2^{-k}$, and every $k$ occurs. Writing
$M_K=\max_{x\in K}|x|$, equation~\eqref{iterate} gives the exact formula
\begin{equation}\label{distance}
 h(F_t^n(K),\{0\})=r_t^n M_K.
\end{equation}
Consequently $A_t=\{0\}$ is the attractor for every $t<1$, with basin
$\R$. In particular every increasing sequence of parameters tending
to one has the same Hausdorff limit $\{0\}$.

At the endpoint, $F_1=\{\operatorname{id},x\mapsto x/2\}$ and
$F_1(\{0\})=\{0\}$. This is the smallest invariant set containing
$Q$, so Definition 4.1 gives $\Al=\{0\}$. Definition 4.3 gives the
unique upper transition attractor $\Au=\{0\}$.

Let $U$ be any open neighborhood of zero. It contains a nonzero point
$x$, and $K=\{x\}\subset U$ is compact. Formula~\eqref{distance}, now
with $r_1=1$, says
\[
 h(F_1^n(\{x\}),\{0\})=|x|>0\qquad(n\geq0).
\]
Thus no such $U$ is a basin of attraction for $\{0\}$. This disproves
the conjecture's conjunction, even though its uniqueness conclusion
holds in this example.
\end{proof}

The source's finite-contractivity-threshold convention also causes
no restriction: the same formulas extend to all $t\geq0$. For $t>1$,
the orbit $g_t^n(x)=r_t^nx$ of any nonzero $x$ is unbounded, ruling out
an attractor with an open basin. It also rules out contraction in an
equivalent complete metric, since such a contraction would make every
$g_t$ orbit converge to its fixed point zero. At $t=1$ the identity
cannot be a contraction in any metric. Both thresholds are therefore
one.

\paragraph{Task 2. Characterize singleton attraction for linear IFSs.}
Express strict attraction to a common fixed point in terms of uniform
decay of all matrix products, and identify the role of the joint
spectral radius in finite dimensions.

\paragraph{Task 3. Settle the separate uniqueness implication.}
Determine whether $\Al=\Au$ for one subsequential limit forces all
upper transition attractors to agree under (H1)--(H3).

\paragraph{Task 4. Formulate a sufficient terminal-attraction condition.}
Find a quantitative condition on finite iterates near the common
transition set that guarantees an open basin at $t=1$.

\paragraph{Task 5. Construct examples with larger transition sets.}
Extend the failure of terminal attraction to families whose coinciding
transition attractors have positive dimension, and compare periodic
and nonperiodic endpoint isometries.

\begin{thebibliography}{1}
\bibitem{source} K. Le\'sniak, N. Snigireva, F. Strobin and A. Vince,
\emph{Transition phenomena for the attractor of an iterated function
system}, Nonlinearity \textbf{35} (2022), 5396--5426.
\href{https://doi.org/10.1088/1361-6544/ac8af1}{doi:10.1088/1361-6544/ac8af1}.
\end{thebibliography}
\end{document}
